\documentclass[11pt,a4paper]{article}

% -------------------------------------------------------------------------
%  PACKAGES  (preamble inherited from proposal.tex so the two documents
%  sit together as one body of work)
% -------------------------------------------------------------------------
\usepackage[T1]{fontenc}
\usepackage[utf8]{inputenc}
\usepackage{mathpazo}
\usepackage{amsmath,amssymb}
\usepackage{geometry}
\usepackage{xcolor}
\usepackage{titlesec}
\usepackage{enumitem}
\usepackage{booktabs}
\usepackage{array}
\usepackage{multirow}
\usepackage{tabularx}
\usepackage{tcolorbox}
\tcbuselibrary{skins,breakable}
\usepackage{fancyhdr}
\usepackage{graphicx}
\usepackage{hyperref}
\usepackage{caption}
\usepackage{placeins}
\usepackage{tikz}
\usetikzlibrary{arrows.meta,positioning}

\geometry{a4paper,left=2.6cm,right=2.6cm,top=2.6cm,bottom=2.8cm}
\setlength{\headheight}{22pt}

% -------------------------------------------------------------------------
%  PALETTE
% -------------------------------------------------------------------------
\definecolor{Ink}{HTML}{203247}
\definecolor{Accent}{HTML}{2E6E8E}
\definecolor{Highlight}{HTML}{C5681B}
\definecolor{AccentPale}{HTML}{EAF2F5}
\definecolor{Paper}{HTML}{FBFBF9}
\definecolor{Muted}{HTML}{5F6F7B}
\definecolor{Rule}{HTML}{C9D6DC}
\definecolor{Warn}{HTML}{8C3B2E}

\pagecolor{Paper}
\color{Ink}

\hypersetup{
  colorlinks=true,
  linkcolor=Accent,
  citecolor=Accent,
  urlcolor=Accent,
  pdftitle={Rooting for Kepler: Final Report},
  pdfauthor={CSE 402 Subsection A2}
}

% -------------------------------------------------------------------------
%  SECTION TITLING
% -------------------------------------------------------------------------
\titleformat{\section}
  {\normalfont\Large\scshape\color{Ink}}
  {\textcolor{Accent}{\thesection}}{0.6em}{}
  [\vspace{-0.3em}{\color{Rule}\titlerule[0.8pt]}]
\titlespacing*{\section}{0pt}{1.6em}{0.9em}

\titleformat{\subsection}
  {\normalfont\large\scshape\color{Accent}}
  {\thesubsection}{0.55em}{}
\titlespacing*{\subsection}{0pt}{1.1em}{0.4em}

% -------------------------------------------------------------------------
%  HEADER / FOOTER
% -------------------------------------------------------------------------
\pagestyle{fancy}
\fancyhf{}
\renewcommand{\headrulewidth}{0.5pt}
\renewcommand{\footrulewidth}{0pt}
\fancyhead[L]{\footnotesize\scshape\color{Muted}CSE 402 \textbar\ Numerical Methods}
\fancyhead[R]{\footnotesize\scshape\color{Muted}Final Report}
\fancyfoot[C]{\footnotesize\color{Muted}\thepage}
\renewcommand{\headrule}{\hbox to\headwidth{\color{Rule}\leaders\hrule height \headrulewidth\hfill}}

% -------------------------------------------------------------------------
%  UTILITIES
% -------------------------------------------------------------------------
\newcommand{\key}[1]{\textcolor{Highlight}{\bfseries #1}}

\newtcolorbox{note}{
  enhanced, breakable, colback=AccentPale, colframe=Rule,
  boxrule=0.5pt, arc=0pt, outer arc=0pt,
  left=8pt, right=8pt, top=6pt, bottom=6pt,
  before skip=1em, after skip=1em
}
\newtcolorbox{notetitled}[1]{
  enhanced, breakable, colback=AccentPale, colframe=Rule,
  boxrule=0.5pt, arc=0pt, outer arc=0pt,
  left=8pt, right=8pt, top=6pt, bottom=6pt,
  before skip=1em, after skip=1em,
  fonttitle=\scshape\footnotesize\color{Accent},
  colbacktitle=AccentPale, coltitle=Accent,
  title={#1}
}
\newtcolorbox{finding}[1]{
  enhanced, breakable, colback=Paper, colframe=Accent,
  boxrule=0.9pt, arc=0pt, outer arc=0pt,
  left=9pt, right=9pt, top=7pt, bottom=7pt,
  before skip=1em, after skip=1.1em,
  fonttitle=\scshape\footnotesize\color{Accent},
  colbacktitle=Paper, coltitle=Accent,
  title={#1}
}

\newcommand{\dividerline}{\vspace{0.2em}{\color{Rule}\hrule height 0.5pt}\vspace{0.6em}}

\renewcommand{\arraystretch}{1.25}

\setlist[itemize]{label=\textcolor{Accent}{\small$\bullet$}, leftmargin=1.4em, itemsep=0.15em}
\setlist[enumerate]{leftmargin=1.6em, itemsep=0.2em}

% Figure helper: the figures/ directory is git-ignored and regenerated by
% scripts/make_report_figures.py, so the document must still compile for a
% reader who has not run the pipeline.
\newcommand{\pipelinefigure}[3]{%
  \begin{figure}[htbp]\centering
  \IfFileExists{#1}{\includegraphics[width=#2\textwidth]{#1}}{%
    \fbox{\parbox{0.8\textwidth}{\footnotesize\color{Muted}\centering
      \texttt{#1} not present.\\ Run \texttt{python scripts/make\_report\_figures.py}
      to regenerate every figure from \texttt{results/}.}}}
  \caption{#3}
  \end{figure}}

% =========================================================================
\begin{document}
\pagenumbering{gobble}

\begin{titlepage}
\pagecolor{Paper}
\thispagestyle{empty}
\begin{center}
  \vspace*{1.6cm}
  {\color{Accent}\scshape\large CSE 402 \,\textbar\, Numerical Methods\, --\, Final Report}\\[0.6cm]
  {\color{Rule}\rule{0.55\textwidth}{0.8pt}}\\[1.1cm]

  {\color{Ink}\scshape\Huge Rooting for Kepler}\\[0.35cm]
  {\color{Muted}\Large\itshape a numerical benchmark of\\[2pt] Kepler solvers and starting guesses}\\[1.4cm]

  {\color{Rule}\rule{0.32\textwidth}{0.6pt}}\\[1.1cm]

  \begin{minipage}{0.76\textwidth}
    \centering
    \small\color{Muted}
    An empirical study of classical, domain-specialised, and 2024--2025
    with-memory Newton-type variants applied to Kepler's equation,
    evaluated on RadVel's radial-velocity fitting pipeline.\\[0.7em]
    \normalcolor\color{Ink}
    125{,}800 instrumented solves over a 7{,}400-point $(e,M)$ grid,
    with iterative solvers checked against published convergence orders
    before entering the benchmark.
  \end{minipage}

  \vspace{2.0cm}
  {\scshape\color{Accent}\footnotesize Subsection A2 \quad\textbar\quad Student IDs}\\[0.35cm]
  {\color{Ink}2105036 \quad 2105050 \quad 2105053 \quad 2105056 \quad 2105059}

  \vfill
  {\footnotesize\color{Muted}\today}
\end{center}
\end{titlepage}

\pagenumbering{arabic}

% =========================================================================
\section*{Summary of Findings}
\addcontentsline{toc}{section}{Summary of Findings}

The project compares recent with-memory root finders with established
Kepler solvers, varying both the starting guess and the orbital regime.
Higher formal order reduces the number of iterations, but does not by
itself imply a cheaper or more reliable solve.

\begin{enumerate}
  \item \key{The verification gates pass.} On eight test functions from
        the NWM11 paper, measured orders agree with the published values
        within $1\%$; the two with-memory base error constants also match
        their respective papers. NWM9 has not yet been tested on its own
        paper's example functions.
  \item \key{Fewer iterations do not ensure lower cost.} With
        selected starting guesses, NWM11 needs a median of one iteration
        against Danby's two. Its modelled median cost is $6.34$ versus
        $7.68$, but its \emph{mean cost among converged solves} is
        $9.55$ versus $7.48$.
  \item \key{The timed Python implementations favour simpler methods.}
        On the timed subset, the median solve takes $5.25\,\mu\mathrm{s}$ for
        NWM11 with its canonical start and $4.46\,\mu\mathrm{s}$ for Danby
        with Napier's start. These are implementation timings, not a
        compiled-algorithm comparison.
  \item \key{Reliability depends on the starting guess.} Across all four
        guesses, NWM9 and NWM11 fail on $4.27\%$ and $4.86\%$ of solves in
        the defined hard corner, versus Danby's $0.13\%$. NWM9 with
        Napier's start has no failures on this grid.
  \item \key{The proposed fused-trigonometry benefit remains untested.}
        This Python code evaluates $\sin$ and $\cos$ separately. Under
        one-pair-per-call accounting, the representative NWM11 and Danby
        solves both cost five calls; a compiled speed advantage cannot be
        inferred from these data.
  \item \key{Downstream shifts are small in the tested injection.}
        Fitted-parameter shifts from varying solver tolerance are at least
        about $10^3$ times smaller than the Monte Carlo spread for this
        moderately eccentric injected orbit.
\end{enumerate}

\begin{finding}{Practical reading of the benchmark}
Markley is fastest on the timed grid subset and has no failures there;
its closed-form computation is not fully captured by the trigonometric
cost proxy. Danby remains a strong iterative baseline. The with-memory
schemes reduce iteration counts, but this experiment does not establish a
total-cost or wall-clock advantage over Danby.
\end{finding}

% =========================================================================
\section{What Was Built}

The deliverable is a cost-instrumented benchmark harness
(\texttt{keplerbench}) with a shared solver interface and configured
experiments. The results and figures were regenerated using the commands
below; the numerical tables in this document were then checked against
those outputs.

\subsection{Architecture}

All five solvers and four starting guesses expose the same calling
interface. The four iterative solvers use a residual stopping test;
Markley is a closed-form method that ignores the requested tolerance and
starting guess. Solver evaluations pass through a
\texttt{KeplerProblem} object that records what each call needed
--- how many distinct points were touched, how many required a full
$(\sin,\cos)$ pair against $\sin$ or $\cos$ alone, and how many
derivatives were synthesised from interpolation rather than evaluated.
That instrumentation is what makes Section~\ref{sec:cost} possible.

Reference roots come from an independent 50-digit bracketing solve in
\texttt{mpmath}, never from one of the methods under test.

\subsection{Reproducibility}

The benchmark and verification raw tables have provenance sidecars with
their configuration and source commit. Several derived and downstream
CSVs do not yet have sidecars, so the blanket provenance claim would be
incorrect. The results used here were regenerated from implementation
commit \texttt{48dcad1} with Python 3.12.6; the configured cost weights
were calibrated separately on Python 3.14.7.

\begin{center}
\small
\begin{tabularx}{\textwidth}{@{}>{\raggedright\arraybackslash}p{6.1cm} >{\raggedright\arraybackslash}X@{}}
\toprule
\textbf{Command} & \textbf{Produces} \\
\midrule
\texttt{scripts/run\_verification.py} & \texttt{results/verification/} --- order and correctness gate \\
\texttt{scripts/run\_grid\_benchmark.py} & \texttt{results/grid\_benchmark/} --- 125{,}800 solves \\
\texttt{scripts/run\_error\_propagation.py} & \texttt{results/error\_propagation/} --- tolerance sweep \\
\texttt{scripts/run\_monte\_carlo.py} & \texttt{results/monte\_carlo/} --- noise realisations \\
\texttt{scripts/make\_report\_figures.py} & \texttt{figures/} --- every figure below \\
\bottomrule
\end{tabularx}
\end{center}

% =========================================================================
\section{Candidate Methods as Implemented}
\label{sec:methods}

\begingroup
\small
\begin{tabularx}{\textwidth}{@{}>{\raggedright\arraybackslash}p{1.9cm} >{\raggedright\arraybackslash}p{3.1cm} >{\raggedright\arraybackslash}p{1.25cm} >{\raggedright\arraybackslash}X@{}}
\toprule
\textbf{Category} & \textbf{Method} & \textbf{Order} & \textbf{Role in the benchmark} \\
\midrule
Baseline & Newton--Raphson & 2 & Textbook anchor; one evaluation point, one $(\sin,\cos)$ pair per iteration. \\[0.4em]
Kepler solver & Danby's method & 4 & RadVel's production solver; domain-tuned to this equation. \\
 & Markley's method & \textemdash & Closed-form solver; no iteration, constant cost. \\[0.4em]
With-memory & NWM9 (2024) & 8.8989 & Single self-accelerating parameter in the \emph{third} step of an optimal eighth-order base. \\
 & NWM11 (2025) & 10.7446 & Bi-parametric; 3 evaluation points, 1 full $(\sin,\cos)$ pair per iteration. \\[0.4em]
Guess layer & Napier's guess (2024) & n/a & Symbolic-regression start, run in front of every iterative solver. \\
\bottomrule
\end{tabularx}
\endgroup

\vspace{0.5em}
\begin{notetitled}{Three corrections to the proposal's method table}
\begin{enumerate}[leftmargin=1.4em, itemsep=0.35em]
  \item The with-memory schemes are \key{not derivative-free}. Both use
        $\varphi'(s_k)$ explicitly in their first sub-step. For Kepler this
        is favourable rather than otherwise: $f$ and $f'$ come from a
        single $(\sin,\cos)$ pair, which is exactly why the ``3 points,
        1 pair'' economy is achievable at all.
  \item NWM9 is \key{not the controlled sibling of NWM11}. They come from
        two different papers with two different eighth-order base methods
        and the parameter placed in different sub-steps. NWM11's true
        same-base control is \key{NWM10} (R-order 10), defined in the same
        paper as NWM11. This matters for interpreting
        Section~\ref{sec:robust}, where the two behave qualitatively
        differently.
  \item The phrase ``at no additional evaluation points per iteration''
        is misleading. Both schemes use three evaluation points per
        iteration against Newton's one. What costs nothing extra is the
        \emph{derivative information}, obtained by Hermite interpolation
        over points already paid for.
\end{enumerate}
\end{notetitled}

% =========================================================================
\section{Verification: the Gate Before the Benchmark}
\label{sec:verify}

The iterative solvers were checked for convergence order before the grid
benchmark. This reduces the risk of transcription errors, but is not an
independent proof of each implementation.

\subsection{Measured convergence order}

Each solver was run on the eight nonlinear test functions of the NWM11
paper, from the paper's own starting points, at 2000-digit working
precision. NWM9 comes from a different paper with its own examples, but
it was gated on the same NWM11 set; checking it on its own paper's
functions remains to be done (\texttt{experiments/verification.py}).

\begin{center}
\small
\begin{tabular}{@{}lrrrc@{}}
\toprule
\textbf{Solver} & \textbf{Claimed} & \textbf{Measured (mean)} & \textbf{Range over 8 functions} & \textbf{Gate} \\
\midrule
Newton  & 2.0000  & 2.0000  & 2.0000 -- 2.0000   & PASS \\
Danby   & 4.0000  & 4.0000  & 4.0000 -- 4.0000   & PASS \\
NWM9    & 8.8989  & 8.9114  & 8.8746 -- 8.9816   & PASS \\
NWM11   & 10.7446 & 10.7605 & 10.7147 -- 10.8336 & PASS \\
\bottomrule
\end{tabular}
\end{center}

Every measurement sits slightly \emph{above} the claim, which is correct:
both papers state their R-order as a lower bound (``at least'').

\begin{notetitled}{Why 2000 digits, and not 16}
An order-$p$ method multiplies its correct digits by $p$ each step. From
an error of $10^{-1}$, an order-9 method passes $10^{-9}$, $10^{-81}$,
$10^{-729}$. The estimator needs three consecutive usable errors to form
one estimate, so the working precision must outrun $p^3$ digits or the
sequence hits the floor before a single estimate exists. At 100 digits
NWM9 measured anywhere in $8.87$--$10.19$; at 2000 it settles to
$8.87$--$8.98$, and 3000 buys nothing further. Double precision cannot
measure these methods at all.
\end{notetitled}

\subsection{The order alone is not sufficient evidence}

Both papers' iteration formulas contain nested fractions whose PDF text
layer admits several readings, and for both schemes \emph{more than one
reading yields the correct order}. Measuring order 8 would therefore not
have detected a mis-transcription.

What does detect it is each paper's explicit error \emph{constant}. Each
scheme was driven for exactly one step from a point $10^{-24}$ from the
root and its eighth-order constant compared against the published closed
form:

\begin{center}
\small
\begin{tabular}{@{}llr@{}}
\toprule
\textbf{Scheme} & \textbf{Published constant} & \textbf{Agreement (relative)} \\
\midrule
NWM9  & Eq.~(7): $c_2^2 K[(\alpha+c_2)K + c_4]$, $K = c_2+5c_2^2-c_3$ & $\leq 1.8\times10^{-22}$ \\
NWM11 & Eq.~(2.6): $(\alpha+c_2)^2 c_3 (c_2c_3 - c_4)$                & $\leq 3.5\times10^{-22}$ \\
\bottomrule
\end{tabular}
\end{center}

A wrong grouping in NWM11's $w'(t_k)$ missed this constant by roughly 60
orders of magnitude while still converging at order 8. Three of the eight
test functions also had to be corrected: superscripts detach in the PDF
text layer, so $e^{s^2}$ extracts as $e^{s}$ followed by a stray $2$.
These were caught because the published roots did not zero the
mis-transcribed functions.

\subsection{Two-stage verification}

Setting each scheme's self-accelerating parameters to zero must recover
the underlying memoryless method. Both do, measuring exactly $8.0000$ on
all eight functions, which attributes the lift to $8.8989$ and $10.7446$
to the memory rather than to a coincidence in the base method.

\subsection{Correctness: the right root, not just a root}

A solver that jumps to a different branch of Kepler's equation reports a
tiny residual and a converged flag while every downstream quantity is
wrong. Over 4{,}975 $(e,M)$ points spanning the full range:

\begin{center}
\small
\begin{tabular}{@{}lrrr@{}}
\toprule
\textbf{Solver} & \textbf{Converged} & \textbf{Wrong root} & \textbf{Max $|E-E_{\mathrm{ref}}|$} \\
\midrule
Newton & 99.54\% & 0 & $3.38\times10^{-10}$ \\
Danby  & 99.72\% & 0 & $1.65\times10^{-10}$ \\
NWM9   & 98.25\% & 0 & $2.32\times10^{-10}$ \\
NWM11  & 99.56\% & 0 & $2.23\times10^{-10}$ \\
\bottomrule
\end{tabular}
\end{center}

No method converged to a wrong branch anywhere. Non-convergence is
reported honestly and is treated as a robustness result
(Section~\ref{sec:robust}), not as a correctness failure.

% =========================================================================
\section{Kepler-Specific Cost Accounting}
\label{sec:cost}

\subsection{Measured per-iteration cost}

Counted from the instrumented code rather than transcribed from the
papers:

\begin{center}
\small
\begin{tabular}{@{}lrrrrr@{}}
\toprule
\textbf{Method} & \textbf{Eval.\ points} & \textbf{$(\sin,\cos)$ pairs} & \textbf{$\sin$ only} & \textbf{Synth.\ derivatives} & \textbf{Order} \\
\midrule
Newton  & 1 & 1 & 1 & 0 & 2.0000 \\
Danby   & 1 & 1 & 1 & 0 & 4.0000 \\
Markley & 2 & 1 & 1 & 0 & --- \\
NWM9    & 3 & 1 & 3 & 3 & 8.8989 \\
NWM11   & 3 & 1 & 3 & 3 & 10.7446 \\
\bottomrule
\end{tabular}
\end{center}

In the measured steady-state iteration, both with-memory schemes use
three evaluation points but only one full $(\sin,\cos)$ pair. Three
derivatives are synthesised by interpolation; the first iteration has
not yet accumulated this information.

\begin{note}
\textbf{One count in this table is an artefact of the harness, not of the
methods.} The shared iteration loop calls $f(E)$ itself after each step to
compute the residual for the stopping test. That inflates the
$\sin$-only column by exactly one per iteration for \emph{every} solver
--- which is why Newton and Danby show a $\sin$-only cost of 1 despite
needing nothing beyond their single $(\sin,\cos)$ pair. The with-memory
schemes' own $\sin$-only cost is therefore 2, not 3. The inflation is
uniform, so it does not distort comparisons, but it must be subtracted
before any of these figures is quoted against a paper's table.
\end{note}

\subsection{Costing assumptions and their limits}

The proposal's argument rests on a specific empirical claim: that one
\texttt{sincos} call returns both $\sin E$ and $\cos E$ ``at roughly the
cost of one transcendental evaluation''. The benchmark uses the following
configured weights, calibrated on a separate CPython 3.14.7 run and
normalised to one $\sin$ call:

\begin{center}
\small
\begin{tabular}{@{}lrr@{}}
\toprule
\textbf{Operation} & \textbf{Relative cost} & \textbf{Proposal assumed} \\
\midrule
$\sin$ only                      & 1.000 & 1 \\
$\cos$ only                      & 1.04 & 1 \\
$(\sin,\cos)$ pair               & \textbf{2.34} & $\approx 1$ \\
Synthesised derivative (Hermite) & 1.76 & --- \\
\bottomrule
\end{tabular}
\end{center}

These are the code's default weights; for example,
Newton's modelled steady-state cost is $2.34+1=3.34$. The current
Python 3.12.6 environment did not reproduce a stable derivative weight:
the default calibration raised an error when its baseline subtraction
gave a non-positive $\sin$ cost, and a longer retry had a wide spread.
Thus weighted cost is an \emph{assumption-sensitive proxy}, not a
measured runtime prediction on this machine. Starting-guess computation,
division, roots and other arithmetic are not fully included.

\begin{finding}{A fused pair does not establish a compiled speed-up}
The Python implementation obtains $\sin$ and $\cos$ with two separate
calls. For the representative median solves, let $p$ be the cost of a
pair in units of one $\sin$ call. NWM11 costs $4+p$ and Danby $3+2p$:
at the configured $p=2.34$ the model gives $6.34$ versus $7.68$, but at
$p=1$ both cost $5$. Even at $p=1.1$ the gap is only about $2\%$.
Across converged solves on the \emph{whole grid}, NWM11 with its canonical
guess has a mean modelled cost of $9.55$ versus Danby with Napier's
$7.48$; the call-count proxy gives $6.30$ versus $4.88$. A compiled
implementation and a direct timing comparison would be needed before
claiming a speed advantage.
\end{finding}

\pipelinefigure{figures/cost/cost_breakdown.pdf}{1.0}{%
Steady-state per-iteration proxy under configured Python weights
(left: $3.34$ for Newton and Danby, $10.62$ for NWM9 and NWM11) and
one-call-per-evaluation-point counting (right: $2$ and $4$). This is
\emph{per iteration}, not per completed solve; the latter also depends on
the number of iterations and the starting guess. Markley's non-trigonometric
arithmetic is not represented. Both panels include one residual check.}

% =========================================================================
\section{Grid Benchmark}
\label{sec:grid}

Every solver--guess combination was run over a 7{,}400-point $(e,M)$ grid
combining a uniform block, a logarithmically-spaced pathological corner
and a RadVel-realistic sample: 17 combinations $\times$ 7{,}400 points
$=$ \key{125{,}800 instrumented solves}. Markley ignores the starting
guess and is run once under the label \texttt{n/a}.

\subsection{Headline comparison}

Each iterative solver is shown with a start that gave strong convergence
on this grid. Mean iterations and modelled cost use converged grid solves;
the wall-clock column uses medians over 50 timed points, each repeated
1000 times. The two cost columns answer different questions and should
not be read as interchangeable.

\begin{center}
\small
\begin{tabular}{@{}llrrrr@{}}
\toprule
\textbf{Solver} & \textbf{Guess} & \textbf{Mean iter.} & \textbf{Median cost} & \textbf{Median wall (µs)} & \textbf{Converged} \\
\midrule
Markley & n/a       & 0.00 & \textbf{3.34} & \textbf{3.50} & 100.00\% \\
NWM11   & canonical & \textbf{1.38} & 6.34 & 5.25 & 99.93\% \\
NWM9    & napier    & 1.44 & 6.34 & 6.52 & 100.00\% \\
Danby   & napier    & 1.94 & 7.68 & 4.46 & 100.00\% \\
Newton  & napier    & 3.21 & 11.02 & 4.40 & 100.00\% \\
\bottomrule
\end{tabular}
\end{center}

The median does not describe the costly tail.

\key{Iteration count.} NWM11 needs a median of one iteration where Danby
needs two and Newton needs three. This is consistent with the measured
convergence orders, but most one-iteration solves cannot benefit from
recycled information from a previous iteration.

\key{Weighted cost.} The chosen NWM11 median is $6.34$ against Danby's
$7.68$, about $17\%$ lower under the configured weights. Their means
reverse the ordering: $9.55$ against $7.48$. Under one-call-per-point
counting the means are $6.30$ against $4.88$. The median alone therefore
does not establish an overall cost advantage.

\key{Wall-clock time.} On this Python implementation and timed subset,
NWM11 and NWM9 with the displayed guesses take about $1.2$ and $1.5$
times Danby's median time. A weak simple start makes NWM9 much slower
($17.69\,\mu\mathrm{s}$); runtime is not fixed by the solver name alone.

\subsection{The guess layer}

The starting guess can change both iteration count and reliability.
Newton's mean count on converged solves moves from $3.84$ with $E_0=M$
to $3.21$ with Napier's formula, and convergence rises from $99.88\%$ to
$100\%$. NWM9 changes from $99.43\%$ with the simple start to $100\%$
with Napier's. The guess cost itself is excluded from the modelled cost
counters but included in wall-clock timing.

This is why the guess is treated as an independent experimental factor.

\pipelinefigure{figures/grid/guess_effect.pdf}{0.62}{%
Mean iterations saved against $E_0=M$ on the full grid. Napier's guess
has the largest improvement for Newton and NWM9; NWM11 instead has its
lowest mean count with the canonical guess.}

\pipelinefigure{figures/convergence/residual_histories.pdf}{1.0}{%
Residual histories at a typical $(e,M)$ and in the pathological corner.
Both with-memory schemes reach the double-precision floor in two
iterations against Newton's four to five. The shaded band marks the floor
below which a curve measures rounding rather than convergence.}

\pipelinefigure{figures/grid/iterations_heatmaps.pdf}{0.78}{%
NWM9 iterations with the simple and Napier starts on the complete uniform
$50\times100$ grid ($0\leq e\leq0.99$, $0\leq M\leq\pi$), on a shared
colour scale. The simple start has a much larger high-iteration region;
the logarithmically sampled extreme corner is shown separately in the
failure maps.}

\pipelinefigure{figures/cost/cost_vs_accuracy.pdf}{1.0}{%
\emph{Left}: attained accuracy against modelled cost per solve; successful
methods commonly reach the double-precision floor. \emph{Right}: the
generic order-per-call index against median cost per correct digit.
Newton does not lead on this measure: its median is $0.70$, versus $0.42$
for NWM11 and $0.48$ for Danby. The two rankings are broadly similar;
this experiment does not show the large rank reversal anticipated in the
proposal. The digit calculation includes exact-reference hits at the
project's 16-digit ceiling.}

% =========================================================================
\FloatBarrier
\section{Robustness}
\label{sec:robust}

Failure rates split by region. The current analysis code defines the
hard corner as $e>0.9$ and $M<0.1$; boundary samples at exactly $0.9$
or $0.1$ are classified as ordinary. Counts below are \emph{solves}
across four guesses for each iterative method, not distinct grid points.

\begin{center}
\small
\begin{tabular}{@{}lrrrr@{}}
\toprule
\textbf{Solver} & \textbf{Corner solves} & \textbf{Corner failure} & \textbf{Ordinary solves} & \textbf{Ordinary failure} \\
\midrule
Markley & 381  & 0.000\% & 7{,}019  & 0.000\% \\
Danby   & 1{,}524 & 0.131\% & 28{,}076 & 0.000\% \\
Newton  & 1{,}524 & 0.919\% & 28{,}076 & 0.018\% \\
NWM9    & 1{,}524 & 4.265\% & 28{,}076 & 0.057\% \\
NWM11   & 1{,}524 & 4.856\% & 28{,}076 & 0.000\% \\
\bottomrule
\end{tabular}
\end{center}

\begin{finding}{Robustness is a solver--guess result}
NWM11 has no failures in the ordinary region, but 74 failed solves in
the hard corner across its four guesses. Its canonical guess fails on
5 of 381 corner points; its RadVel guess fails on 35. NWM9 with
Napier's guess has no failures anywhere on this grid, despite 81 failed
solves for NWM9 across all guesses.

The small derivative $f'(E)=1-e\cos E$ makes this corner difficult.
Freezing the accelerating parameters can improve some failures, but
that changes both their initial values and the recycling rule. The
current experiment does not isolate memory recycling as the sole cause
of failure; the starting guess is clearly important.
\end{finding}

\pipelinefigure{figures/grid/failure_maps.pdf}{1.0}{%
Where each solver fails, on logarithmic axes in $1-e$ and $M$ so the
pathological corner occupies the top left rather than a sliver. The map
marks a \emph{distinct point} if any guess fails there: Markley $0$,
Danby $2$, Newton $19$, NWM11 $65$, NWM9 $74$. These need not equal the
failed-solve counts in the table. Points at exactly $e=0$ or $M=0$
cannot be drawn on log axes.}

\subsection{Order across the eccentricity range}

Convergence order measured on Kepler's equation itself at $M = 0.3$, at
2000-digit precision --- a different and more demanding question than the
generic test functions of Section~\ref{sec:verify}.

\begin{center}
\small
\begin{tabular}{@{}lrrrrrrrrr@{}}
\toprule
$e$ & 0.1 & 0.3 & 0.5 & 0.7 & 0.9 & 0.95 & 0.99 & 0.995 & 0.999 \\
\midrule
Newton & 2.000 & 2.000 & 2.000 & 2.000 & 2.000 & 2.000 & 2.129 & --- & 1.786 \\
Danby  & 4.000 & 4.000 & 4.000 & 4.000 & 4.000 & 4.000 & 4.000 & 4.000 & 4.000 \\
NWM9   & 8.891 & 8.897 & 8.910 & 8.896 & \textit{2.820} & 8.899 & 4.479 & 0.721 & \textbf{0.365} \\
NWM11  & 11.233 & 10.716 & 10.728 & 10.756 & 10.876 & 10.839 & 10.733 & 10.739 & \textbf{10.744} \\
\bottomrule
\end{tabular}
\end{center}

Along this $M=0.3$ slice, NWM11's measured order remains near its
claimed $10.7446$ up to $e=0.999$; Danby's is 4.000. This test does
not establish the same behaviour as $M\rightarrow0$.

NWM9's estimated order drops beyond $e=0.95$ in this slice. However,
each extreme estimate uses only a few pre-floor errors, the dip at
$e=0.9$ recovers immediately, and no independent implementation of NWM9
has been checked here. These are observations about the estimator and
implementation, not proof that the mathematical scheme changes order.
Newton's missing or low extreme-$e$ estimates have the same resolution
caveat.

\pipelinefigure{figures/convergence/order_vs_eccentricity.pdf}{0.92}{%
Measured order at $M=0.3$. Dotted lines mark published orders; hollow
markers identify estimates with limited resolution. NWM9's extreme-$e$
estimates need independent verification.}

\pipelinefigure{figures/convergence/measured_vs_claimed_order.pdf}{0.86}{%
Measured convergence order against each source paper's claim, on the
NWM11 paper's eight test functions (used for every solver). Error bars show the spread across the eight
functions.}

% =========================================================================
\FloatBarrier
\section{Downstream: Does Any of It Matter?}
\label{sec:downstream}

Solver error was traced through $M \rightarrow E \rightarrow \nu
\rightarrow v_r$ into refitted orbital parameters, with each of our
solvers patched into RadVel in place of its own compiled solver. A call
counter asserts after every fit that the patch actually fired.

\key{Design: injection-recovery on real cadence.} The observation times,
the per-point uncertainties and the three-instrument structure come from
401 HD~164922 measurements. The velocities are
\emph{not}: they are replaced by a known, moderately eccentric orbit
($P = 20.885258$\,d, $e = 0.35$, $\omega = 1.0$\,rad, $K = 5$\,m\,s$^{-1}$)
plus Gaussian noise at the real uncertainties. Fitting HD~164922's own
velocities with a single-planet model would give a near-circular,
optimiser-dependent eccentricity in a one-planet fit (the system has four
planets). The injection instead provides a known answer to recover. The
fit varies $P$, the
time of conjunction, $\sqrt{e}\cos\omega$, $\sqrt{e}\sin\omega$ and
$\log K$ (RadVel's recommended basis), plus one jitter term per
instrument. Before the final fit, a 30-point multi-start over $(e,\omega)$
using RadVel's native solver picks the starting point. The fit recovers
$e\approx0.342$, $\omega\approx0.997$ and $K\approx5.03$, close to the
injected values for all five solvers. Each fit performs about
$1.3\times10^{5}$ Kepler solves.

\subsection{Cost of a full fit}

\begin{center}
\small
\begin{tabular}{@{}lrrr@{}}
\toprule
\textbf{Solver} & \textbf{Median fit (s)} & \textbf{Slowdown vs.\ native} & \textbf{µs per solve} \\
\midrule
Newton  & 0.95 & 17.3$\times$ & 7.12 \\
Danby   & 1.03 & 18.5$\times$ & 7.35 \\
Markley & 1.00 & 22.3$\times$ & 7.68 \\
NWM11   & 1.23 & 27.8$\times$ & 9.24 \\
NWM9    & 1.35 & 30.6$\times$ & 9.69 \\
\bottomrule
\end{tabular}
\end{center}

Medians over the six tolerances (\texttt{results/error\_propagation/timing.csv}).
Only the fit made with the solver under test is timed, not the multi-start
search. ``Native'' is the identical fit, from the same starting point,
using RadVel's compiled solver (about $0.04$--$0.05$\,s in this run).
Each fit was timed once, so small differences are not reliable rankings.

The slowdown against native RadVel includes Python-to-scalar conversion:
our solvers are scalar and RadVel passes arrays. This does not
isolate algorithmic speed from implementation overhead. On this setup,
the with-memory fits take longer than the classical ones.

\pipelinefigure{figures/cost/wallclock_vs_cost.pdf}{0.72}{%
One point per timed $(e,M)$ solve and solver--guess pair ($850$ points).
Pearson $r\approx0.99$ is dominated by five very slow corner solves;
without them it falls to $0.83$, and the full-sample rank correlation is
$0.49$. The proxy does not reliably rank ordinary solves by runtime.}

\subsection{The error budget}

The question is not how large the solver-induced parameter shift is, but
how it compares to the two uncertainties that are already present.

\begin{center}
\small
\begin{tabular}{@{}lrrrr@{}}
\toprule
\textbf{Parameter} & \textbf{Solver shift} & \textbf{Noise spread} & \textbf{MCMC posterior} & \textbf{Noise / solver} \\
 & \footnotesize(worst case) & \footnotesize(Monte Carlo) & \footnotesize($1\sigma$) & \\
\midrule
$P$ (d)  & $1.0\times10^{-8}$ & $5.04\times10^{-4}$ & $6.44\times10^{-4}$ & $5.1\times10^{4}$ \\
$e$      & $1.1\times10^{-5}$ & $1.65\times10^{-2}$ & $1.75\times10^{-2}$ & $1.5\times10^{3}$ \\
$\omega$ (rad) & $2.3\times10^{-5}$ & $5.03\times10^{-2}$ & $5.59\times10^{-2}$ & $2.2\times10^{3}$ \\
$K$ (m\,s$^{-1}$) & $3.5\times10^{-5}$ & $9.93\times10^{-2}$ & $1.04\times10^{-1}$ & $2.8\times10^{3}$ \\
\bottomrule
\end{tabular}
\end{center}

The solver shift is the largest absolute shift over all five solvers and
all tolerances, measured from each solver's own $10^{-14}$ fit. In every
parameter that worst case is Newton at $10^{-4}$. The noise spread is the
standard deviation over 50 refits of fresh noise realisations. The MCMC
width comes from one RadVel run with the native solver, configured with
\texttt{nrun=2000}; RadVel's convergence checks may extend the total run.

\begin{finding}{Numerical precision is not limiting in this injection}
For this $e=0.35$ injected orbit, even the loosest tested tolerance
($10^{-4}$) gives parameter shifts at least about $10^3$ times smaller
than the spread across 50 noisy refits. Newton at $10^{-4}$ has the
largest observed effect, $6.3\times10^{-4}\sigma$ in $e$. At other
settings the shifts are at most about $1.1\times10^{-5}\sigma$; their
non-monotonic variation is compatible with fit/optimiser resolution and
is best treated as an upper bound on solver effects. Markley ignores
the tolerance, so its six fits are identical; the $10^{-14}$ row is zero
by construction for every solver.

The Monte Carlo spread and MCMC posterior width have similar scales in
this example, but they are different distributions: repeated noisy
refits measure estimator sampling variation, whereas MCMC conditions
on the observed data and priors. Agreement is a check of scale, not an
expectation that the two distributions coincide.
\end{finding}

\pipelinefigure{figures/propagation/amplification.pdf}{0.62}{%
The amplification $\mathrm{d}\nu/\mathrm{d}E$ carrying solver error from the
eccentric anomaly into the true anomaly. Near-circular orbits pass error
through roughly unchanged; at $e = 0.99$ the factor exceeds $14$ at the ends
of the range. This illustrates a possible sensitivity; the downstream
fit experiment used $e=0.35$, not $e=0.99$.}

\pipelinefigure{figures/propagation/tolerance_vs_parameter_shift.pdf}{1.0}{%
Shift in each fitted parameter as the solver tolerance is loosened, in units
of the MCMC posterior $\sigma$. The dotted guides mark $1\sigma$ and
$0.1\sigma$. All plotted shifts are below $10^{-3}\sigma$ for this
injected orbit. Hollow markers denote zero shifts: Markley at every
tolerance and every solver at the $10^{-14}$ reference.}

\pipelinefigure{figures/propagation/error_budget.pdf}{0.78}{%
Solver-induced shift, repeated-noise spread and MCMC posterior width for
the injected orbit. The latter two have similar scales but different
statistical meanings; the maximum solver shift is at least three orders
of magnitude smaller than the noise spread.}

\subsection{Real-data checks}

The main study uses injected velocities, so we also fitted all five
solvers to three datasets' \emph{real}, unmodified velocities at
one tight tolerance. This separate check tests agreement between
solvers, not recovery of a known orbit. For the two transiting planets,
period and transit time are held fixed from external measurements:
\begin{itemize}
  \item K2-24~b: $P = 20.885258$\,d, $t_c = 2072.79438$.
  \item K2-131~b: $P = 0.3693038$\,d, $t_c = 2457582.9360$ [9].
\end{itemize}
HD~164922~b does not transit, so only its period ($1207$\,d) is fixed.
The fits use the same basis and starting-point search for every solver.

\begin{center}
\small
\begin{tabular}{@{}lrr@{}}
\toprule
\textbf{Dataset} & \textbf{Points} & \textbf{Fitted $e$ (approximately)} \\
\midrule
K2-131 (1 planet)    & 70  & 0.135 \\
HD~164922 (4 planets) & 401 & 0.122 \\
K2-24 (2 planets)    & 32  & 0.908 \\
\bottomrule
\end{tabular}
\end{center}

The five solver outputs agree to the displayed precision. HD~164922 and
K2-24 are multi-planet systems fitted here with one planet, so the
reported eccentricities should not be interpreted as new measurements
of their orbits. These direct checks were rerun separately from the
configured injection experiment; their 15 fit rows are saved with the
experiment evidence as \texttt{real\_data\_checks.csv}.

% =========================================================================
\FloatBarrier
\section{Synthesis}

The results separate formal convergence, completed-solve cost and the
effect on a fitted orbit:

\begin{center}
\small
\begin{tabularx}{\textwidth}{@{}>{\raggedright\arraybackslash}p{2.4cm} >{\raggedright\arraybackslash}X@{}}
\toprule
\textbf{Axis} & \textbf{Verdict} \\
\midrule
Convergence order & On the generic verification functions, NWM11 measures 10.76 versus Danby's 4.00. Its order on Kepler was also near the published value along the tested $M=0.3$ slice. \\
Iteration count & Selected NWM11 and NWM9 starts have a median of one iteration versus Danby's two. \\
Modelled cost & The selected NWM11 median is $6.34$ versus $7.68$ for Danby, but the corresponding converged-solve means are $9.55$ versus $7.48$; the weights are machine-sensitive. \\
Python wall-clock & Timed-grid medians for those starts are $5.25$ and $4.46\,\mu\mathrm{s}$ per solve. Full-fit medians also favour the simpler implementations on this setup. \\
Robustness & Failures concentrate near $e\rightarrow1$, $M\rightarrow0$ and depend strongly on the guess. NWM9 with Napier has none on this grid. \\
Downstream fit & For one injected $e=0.35$ orbit, the largest parameter shifts are at least about $10^3$ below the repeated-noise spread. \\
\bottomrule
\end{tabularx}
\end{center}

The proposal's numerical-cost hypothesis is therefore only partly
supported: fewer iterations and a lower \emph{selected median} proxy do
not translate into a lower average cost or a Python runtime win. A fused
pair would erase the representative median proxy gap at $p=1$; compiled
speed remains an experiment to perform, not a conclusion of this one.

Markley has the lowest trigonometric proxy ($3.34$) and fastest median
on the timed grid subset, but the proxy omits some of its arithmetic and
its closed-form interface does not enforce the iterative tolerance.
Six iterative solver--guess pairs also have zero failures on this grid.
Danby remains the appropriate production reference among the iterative
methods studied here.

% =========================================================================
\section{Threats to Validity}

\begin{enumerate}
  \item \key{Implementation and calibration.} Wall-clock timing compares
        scalar Python implementations, not compiled kernels. The cost
        proxy uses weights calibrated under Python 3.14.7, whereas this
        rerun used Python 3.12.6; local recalibration was unstable. It
        also omits starting-guess cost and some non-trigonometric work.
  \item \key{The residual-evaluation artefact.} The shared loop's own
        $f(E)$ call inflates every solver's $\sin$-only count by one per
        iteration. It is uniform across methods and does not affect
        some rankings, but the absolute per-iteration figures in
        Section~\ref{sec:cost} are one $\sin$ high. Markley is not subject
        to the same tolerance stopping test as the iterative solvers.
  \item \key{One injected orbit.} The error-propagation and Monte Carlo
        studies use one injected orbit ($e=0.35$) on HD~164922's real
        cadence and uncertainties, not its own velocities. High-$e$
        injections were not run; the real-data checks only establish
        agreement among solvers under the selected fit model.
  \item \key{Optimiser resolution floor.} Below tolerance $10^{-4}$ the
        measured shifts (outside Newton at $10^{-4}$, at most about
        $1.1\times10^{-5}\sigma$) are compatible with optimiser noise;
        they bound rather than directly measure solver effects.
  \item \key{Uncertainty estimates.} The Monte Carlo sample has only 50
        realisations. The MCMC width comes from one RadVel run configured
        with \texttt{nrun=2000}; RadVel may extend the total run to satisfy
        mixing checks. Sampling spread and posterior width need not match.
  \item \key{Region and provenance.} The hard-corner classifier excludes
        its boundary samples. Several derived/downstream CSVs lack
        provenance sidecars, although the benchmark raw data has one.
  \item \key{Estimator artefacts at extreme eccentricity.} Several entries
        in the order-versus-eccentricity table are measurement artefacts
        rather than method behaviour, as noted in
        Section~\ref{sec:robust}. They are reported rather than removed,
        with the unreliable points marked.
  \item \key{NWM9's degradation is single-sourced.} The collapse beyond
        $e \approx 0.99$ has been reproduced across precisions but not
        cross-checked against an independent implementation of the same
        published scheme.
\end{enumerate}

% =========================================================================
\section{Conclusion}

Five Kepler solvers were checked through one instrumented harness over
125{,}800 solves. The iterative implementations passed the project's
order and correctness gates on the selected verification sets, but NWM9
still needs a check on its own paper's examples and an independent
implementation at extreme eccentricity.

The newer with-memory schemes need fewer iterations with suitable
starts. Their selected median proxy cost is lower than Danby's, while
their average proxy cost and measured Python runtime are higher.
Reliability in the hard corner depends strongly on the starting guess;
it cannot be attributed to memory alone from the present evidence.

In the one injected-orbit fit studied here, solver-tolerance changes are
small compared with noise-driven parameter spread. Danby is a sound
iterative production baseline for this implementation. Whether a
compiled version of either newer method is faster remains untested.

% =========================================================================
\clearpage
\section*{References}
\begingroup
\small
\begin{enumerate}[label={[\arabic*]}, leftmargin=2.2em, itemsep=0.7em]

  \item B. J. Fulton, E. A. Petigura, S. Blunt, and E. Sinukoff,
  ``RadVel: The Radial Velocity Modeling Toolkit,''
  \textit{Publications of the Astronomical Society of the Pacific},
  vol. 130, no. 986, p. 044504, 2018.
  doi: \href{https://doi.org/10.1088/1538-3873/aaaaa8}
  {10.1088/1538-3873/aaaaa8}.

  \item J. M. A. Danby,
  ``The solution of Kepler's equation, III,''
  \textit{Celestial Mechanics},
  vol. 40, nos. 3--4, pp. 303--312, 1987.
  doi: \href{https://doi.org/10.1007/BF01235847}
  {10.1007/BF01235847}.

  \item F. L. Markley,
  ``Kepler equation solver,''
  \textit{Celestial Mechanics and Dynamical Astronomy},
  vol. 63, no. 1, pp. 101--111, 1995.
  doi: \href{https://doi.org/10.1007/BF00691917}
  {10.1007/BF00691917}.

  \item S. K. Mittal, S. Panday, and L. J{\"a}ntschi,
  ``Enhanced Ninth-Order Memory-Based Iterative Technique for Efficiently
  Solving Nonlinear Equations,''
  \textit{Mathematics},
  vol. 12, no. 22, p. 3490, 2024.
  doi: \href{https://doi.org/10.3390/math12223490}
  {10.3390/math12223490}.

  \item S. K. Mittal, S. Panday, L. J{\"a}ntschi, and L. C. Bolundu{\c{t}},
  ``Two novel efficient memory-based multi-point iterative methods for
  solving nonlinear equations,''
  \textit{AIMS Mathematics},
  vol. 10, no. 3, pp. 5421--5443, 2025.
  doi: \href{https://doi.org/10.3934/math.2025250}
  {10.3934/math.2025250}.

  \item K. J. Napier,
  ``Improved Initial Guesses for Numerical Solutions of Kepler's
  Equation,''
  \textit{arXiv preprint arXiv:2411.15374}, 2024.
  doi: \href{https://doi.org/10.48550/arXiv.2411.15374}
  {10.48550/arXiv.2411.15374}.

  \item E. A. Petigura \textit{et al.},
  ``Two Transiting Low Density Sub-Saturns from K2,''
  \textit{The Astrophysical Journal}, vol. 818, no. 1, p. 36, 2016.
  doi: \href{https://doi.org/10.3847/0004-637X/818/1/36}
  {10.3847/0004-637X/818/1/36}.
  \textit{(K2-24 system; see \texttt{data/README.md} for the full dataset
  provenance.)}

  \item B. J. Fulton \textit{et al.},
  ``Three Temperate Neptunes Orbiting Nearby Stars,''
  \textit{The Astrophysical Journal}, vol. 830, no. 1, p. 46, 2016.
  doi: \href{https://doi.org/10.3847/0004-637X/830/1/46}
  {10.3847/0004-637X/830/1/46}.
  \textit{(HD~164922 system.)}

  \item F. Dai \textit{et al.},
  ``The Discovery and Mass Measurement of a New Ultra-short-period Planet:
  K2-131b,''
  \textit{The Astronomical Journal}, vol. 154, no. 6, p. 226, 2017.
  doi: \href{https://doi.org/10.3847/1538-3881/aa9065}
  {10.3847/1538-3881/aa9065}.
  \textit{(K2-131 ephemeris used in the real-data checks.)}

\end{enumerate}
\endgroup

\end{document}
